\chapter{آشنایی با داده و پیش‌پردازش}
\label{ch:data}

\section*{اهداف این فصل}
\begin{itemize}
    \item شناخت انواع ویژگی‌ها و مقیاس‌های اندازه‌گیری داده
    \item درک مفهوم کیفیت داده و منابع اصلی مشکلات آن
    \item تسلط بر روش‌های پاک‌سازی، یکپارچه‌سازی، تبدیل و کاهش داده
    \item آشنایی با معیارهای شباهت و فاصله میان نمونه‌ها
    \item پیاده‌سازی عملی هر مرحله با پایتون و مشاهدهٔ اثر آن روی داده
\end{itemize}

\section{انواع ویژگی‌ها}
\label{sec:attribute-types}

هر مجموعه داده از تعدادی \dmterm{\textbf{نمونه}}{Instance / Object / Record}
تشکیل شده که هر یک با مجموعه‌ای از \dmterm{\textbf{ویژگی‌ها}}{Attributes / Features}
توصیف می‌شود. شناخت دقیق نوع هر ویژگی، پیش‌نیاز انتخاب روش پیش‌پردازش و الگوریتم
مناسب است. ویژگی‌ها را می‌توان از دو منظر دسته‌بندی کرد.

\subsection{دسته‌بندی کیفی در برابر کمّی}

\begin{itemize}
    \item \dmterm{\textbf{ویژگی‌های کیفی}}{Categorical / Qualitative}
    \begin{itemize}
        \item \dmterm{\textit{اسمی}}{Nominal}: مقادیر فقط برچسب‌اند و ترتیب یا فاصلهٔ
              معناداری ندارند. مثال: رنگ چشم، شهر محل سکونت.
        \item \dmterm{\textit{ترتیبی}}{Ordinal}: مقادیر ترتیب معناداری دارند، اما فاصلهٔ
              میان آن‌ها مشخص یا یکسان نیست. مثال: سطح تحصیلات (دیپلم، کارشناسی،
              کارشناسی‌ارشد)، رضایت مشتری (کم، متوسط، زیاد).
    \end{itemize}
    \item \dmterm{\textbf{ویژگی‌های کمّی}}{Numeric / Quantitative}
    \begin{itemize}
        \item \dmterm{\textit{بازه‌ای}}{Interval}: فاصله میان مقادیر معنادار است، اما
              صفر مطلق ندارد. مثال: دما بر حسب سلسیوس.
        \item \dmterm{\textit{نسبتی}}{Ratio}: علاوه بر فاصلهٔ معنادار، صفر مطلق نیز دارد،
              بنابراین نسبت دو مقدار هم معنادار است. مثال: سن، درآمد، وزن.
    \end{itemize}
\end{itemize}

هر نوع ویژگی مجموعهٔ مشخصی از عملیات ریاضی معنادار را مجاز می‌کند؛ برای مثال
میانگین‌گرفتن از «رنگ چشم» بی‌معناست، اما برای «سن» کاملاً معنادار است. جدول
\ref{tab:attr-types} این ارتباط را خلاصه می‌کند.

\begin{table}[htbp]
\centering
\caption{انواع ویژگی‌ها و عملیات معنادار برای هر نوع}
\label{tab:attr-types}
\begin{tabular}{@{}p{2.0cm}p{3.9cm}p{3.4cm}p{2.7cm}@{}}
\toprule
\textbf{نوع} & \textbf{ویژگی کلیدی} & \textbf{مثال} & \textbf{عملیات مجاز} \\
\midrule
اسمی & فقط برچسب، بدون ترتیب & رنگ چشم، کد شهر & برابری، مد \\
\addlinespace
ترتیبی & دارای ترتیب، بدون فاصلهٔ معنادار & سطح تحصیلات، رتبه & مقایسه، میانه \\
\addlinespace
بازه‌ای & فاصله معنادار، بدون صفر مطلق & دمای سلسیوس، سال تقویمی & جمع، تفریق، میانگین \\
\addlinespace
نسبتی & فاصله معنادار و صفر مطلق & سن، درآمد، وزن & همهٔ عملیات، نسبت \\
\bottomrule
\end{tabular}
\end{table}

\subsection{دسته‌بندی گسسته در برابر پیوسته}

از منظری دیگر، ویژگی‌ها به \textbf{گسسته} (تعداد محدود یا شمارش‌پذیر از
مقادیر، مانند تعداد فرزندان) و \textbf{پیوسته} (بی‌نهایت مقدار ممکن در یک
بازه، مانند قد یا وزن) تقسیم می‌شوند. تشخیص درست این نوع، در انتخاب روش
گسسته‌سازی (بخش \ref{subsec:discretization}) و همچنین نوع نمودار مناسب برای
تحلیل اکتشافی داده اهمیت دارد.

\section{کیفیت داده}
\label{sec:data-quality}

پیش از هر کاوشی، باید کیفیت داده ارزیابی شود. مهم‌ترین مشکلات رایج در داده‌های
واقعی عبارت‌اند از:

\begin{itemize}
    \item \dmterm{\textbf{مقادیر گمشده}}{Missing Values}: به دلایلی مانند خطای ثبت،
          عدم پاسخ‌گویی کاربر یا خرابی حسگر.
    \item \dmterm{\textbf{نویز}}{Noise}: خطای تصادفی یا انحراف در مقدار اندازه‌گیری‌شده.
    \item \dmterm{\textbf{داده‌های ناسازگار}}{Inconsistent Data}: تناقض میان رکوردها،
          مثلاً دو مقدار متفاوت برای سن یک فرد در دو جدول مختلف.
    \item \dmterm{\textbf{داده‌های تکراری}}{Duplicate Data}: ثبت چندبارهٔ یک رکورد.
    \item \dmterm{\textbf{داده‌های پرت}}{Outliers}: مقادیری که به‌طور قابل‌توجهی با
          بقیهٔ داده تفاوت دارند (جزئیات بیشتر در فصل \ref{ch:outliers}).
\end{itemize}

نکتهٔ مهم این است که نویز و داده پرت یکسان نیستند: نویز معمولاً خطای تصادفی
و بی‌معنی است که باید حذف یا اصلاح شود، اما داده پرت می‌تواند خودِ هدف تحلیل
باشد (مثلاً تراکنش متقلبانه).

اولین گام عملی در ارزیابی کیفیت، \textbf{دیدن} الگوی مقادیر گمشده است. شکل
\ref{fig:missing-values} نمونه‌ای از این نمودار را نشان می‌دهد: هر سطر یک رکورد،
هر ستون یک ویژگی، و خانه‌های سفید نشان‌دهندهٔ مقدار گمشده‌اند. در این مثال،
ستون‌های «درآمد» و «امتیاز رضایت» بیشترین گمشدگی را دارند و ستون «شهر» تقریباً
کامل است؛ چنین الگویی به تحلیلگر می‌گوید تلاش پاک‌سازی را روی کدام ستون‌ها
متمرکز کند.

\begin{figure}[htbp]
\centering
\includegraphics[width=0.62\linewidth]{figures/ch02_missing_values.pdf}
\caption{نمایش الگوی مقادیر گمشده در یک مجموعه داده (سفید = گمشده)}
\label{fig:missing-values}
\end{figure}

\section{پیش‌پردازش داده}
\label{sec:preprocessing}

پیش‌پردازش داده مجموعه‌ای از گام‌های پشت‌سرهم است که داده خام را به داده‌ای
مناسب برای الگوریتم‌های داده‌کاوی تبدیل می‌کند. شکل \ref{fig:preprocessing-pipeline}
ترتیب متداول این گام‌ها را نشان می‌دهد؛ هر یک در ادامه به‌تفصیل بررسی می‌شود.

\begin{figure}[htbp]
\centering
\begin{tikzpicture}[
    node distance=0.28cm,
    st/.style={rectangle, rounded corners, draw=blue!60, thick, fill=blue!8,
               minimum width=2.0cm, minimum height=0.95cm, text width=1.85cm,
               align=center, font=\scriptsize},
    endp/.style={rectangle, rounded corners, draw=gray!70, thick, fill=gray!12,
               minimum width=2.0cm, minimum height=0.95cm, text width=1.85cm,
               align=center, font=\scriptsize},
    note/.style={text width=2.0cm, align=center, font=\tiny, text=gray!70!black},
    arr/.style={-{Latex[length=2mm]}, thick, draw=blue!60}
]
\node[endp] (raw) {داده خام};
\node[st, left=of raw] (clean) {پاک‌سازی};
\node[st, left=of clean] (integ) {یکپارچه‌سازی};
\node[st, left=of integ] (trans) {تبدیل};
\node[st, left=of trans] (red) {کاهش};
\node[endp, left=of red, fill=green!10, draw=green!60!black] (ready) {داده آماده};

\node[note, below=0.15cm of clean] {حذف نویز و مقدار گمشده};
\node[note, below=0.15cm of integ] {ترکیب چند منبع داده};
\node[note, below=0.15cm of trans] {نرمال‌سازی و گسسته‌سازی};
\node[note, below=0.15cm of red] {نمونه‌برداری و \LR{PCA}};

\draw[arr] (raw) -- (clean);
\draw[arr] (clean) -- (integ);
\draw[arr] (integ) -- (trans);
\draw[arr] (trans) -- (red);
\draw[arr] (red) -- (ready);
\end{tikzpicture}
\caption{گام‌های متداول پیش‌پردازش داده}
\label{fig:preprocessing-pipeline}
\end{figure}

\subsection{پاک‌سازی داده}

رایج‌ترین راهکارها برای مقادیر گمشده در جدول \ref{tab:missing-strategies}
آمده‌اند. انتخاب میان آن‌ها به مقدار گمشدگی، علت آن و اهمیت ستون بستگی دارد.

\begin{table}[htbp]
\centering
\caption{راهکارهای مواجهه با مقادیر گمشده}
\label{tab:missing-strategies}
\begin{tabular}{@{}p{2.9cm}p{4.5cm}p{4.3cm}@{}}
\toprule
\textbf{راهکار} & \textbf{مزیت} & \textbf{عیب} \\
\midrule
حذف رکورد & ساده و سریع & اتلاف داده؛ در گمشدگی زیاد خطرناک \\
\addlinespace
جایگذاری با میانگین/میانه/مد & حفظ تعداد نمونه‌ها؛ پیاده‌سازی آسان & کاهش واریانس و نادیده‌گرفتن روابط میان ویژگی‌ها \\
\addlinespace
جایگذاری با مدل پیش‌بینی & استفاده از سایر ویژگی‌ها؛ دقت بالاتر & پیچیده‌تر و مستعد نشت اطلاعات \\
\bottomrule
\end{tabular}
\end{table}

برای هموارسازی نویز نیز روش‌هایی مانند \dmterm{\textbf{بین‌بندی}}{Binning}،
رگرسیون و تحلیل خوشه‌ای (برای شناسایی و حذف نقاط دورافتاده) به کار می‌روند.

نمونه‌کد زیر، شمارش مقادیر گمشده و جایگذاری آن‌ها با میانه را با
\texttt{pandas} و \texttt{scikit-learn} نشان می‌دهد:

\begin{lstlisting}
import pandas as pd
from sklearn.impute import SimpleImputer

df = pd.read_csv("customers.csv")

# 1) how many values are missing in each column?
print(df.isna().sum())

# 2) fill numeric gaps with the column median
num_cols = df.select_dtypes(include="number").columns
imputer = SimpleImputer(strategy="median")
df[num_cols] = imputer.fit_transform(df[num_cols])
\end{lstlisting}

\subsection{یکپارچه‌سازی داده}

وقتی داده از چند منبع مختلف (مثلاً چند پایگاه داده یا چند فایل) جمع‌آوری
می‌شود، باید به مسائلی مانند موارد زیر توجه کرد:

\begin{itemize}
    \item \dmterm{\textbf{تشخیص موجودیت یکسان}}{Entity Identification}: آیا
          \texttt{customer\_id} در یک جدول همان \texttt{cust\_no} در جدول
          دیگر است؟
    \item \dmterm{\textbf{افزونگی}}{Redundancy}: آیا یک ویژگی از ترکیب چند ویژگی
          دیگر قابل استخراج است؟ (مثلاً سود، از تفاضل درآمد و هزینه.)
    \item \textbf{تناقض مقیاس یا واحد}: مثلاً دما در یک منبع بر حسب سلسیوس و
          در منبع دیگر بر حسب فارنهایت ثبت شده باشد.
\end{itemize}

\subsection{تبدیل و نرمال‌سازی داده}

بسیاری از الگوریتم‌ها (به‌ویژه آن‌هایی که بر پایهٔ فاصله عمل می‌کنند، مانند
نزدیک‌ترین همسایه یا \LR{k-means}) به مقیاس ویژگی‌ها حساس‌اند. اگر یک ویژگی در
بازهٔ $[0, 1]$ و ویژگی دیگر در بازهٔ $[0, 100000]$ باشد، ویژگی دوم به‌طور
نادرست بر محاسبهٔ فاصله غالب می‌شود. دو روش رایج نرمال‌سازی عبارت‌اند از:

\paragraph{نرمال‌سازی \LR{Min-Max}}
مقدار هر ویژگی به بازهٔ $[0, 1]$ نگاشت می‌شود:
\[
x' = \frac{x - \min(x)}{\max(x) - \min(x)}
\]

\paragraph{استانداردسازی \LR{Z-score}}
مقدار هر ویژگی بر اساس میانگین ($\mu$) و انحراف معیار ($\sigma$) آن
تبدیل می‌شود:
\[
x' = \frac{x - \mu}{\sigma}
\]

استانداردسازی \LR{Z-score} معمولاً در حضور داده‌های پرت پایدارتر از
\LR{Min-Max} عمل می‌کند، زیرا \LR{Min-Max} به شدت به مقادیر حداقل و حداکثر
حساس است.

\paragraph{یک مثال عددی}
پنج مشاهده با مقادیر $10, 20, 30, 40, 100$ را در نظر بگیرید (مقدار ۱۰۰ یک
مقدار پرت نسبت به بقیه است). برای این داده $\min = 10$، $\max = 100$،
$\mu = 40$ و $\sigma \approx 31.62$ (انحراف معیار جامعه) است. نتیجهٔ دو
تبدیل در جدول \ref{tab:scaling-example} آمده است. دقت کنید که در \LR{Min-Max}
مقدار پرت باعث می‌شود چهار مقدار دیگر در بازهٔ باریک $[0, 0.33]$ فشرده شوند.

\begin{table}[htbp]
\centering
\caption{مقایسهٔ عددی نرمال‌سازی \LR{Min-Max} و \LR{Z-score} روی یک نمونهٔ کوچک}
\label{tab:scaling-example}
\begin{tabular}{@{}cccc@{}}
\toprule
\textbf{ردیف} & $x$ & \textbf{\LR{Min-Max}} & \textbf{\LR{Z-score}} \\
\midrule
۱ & 10  & 0.000 & $-0.949$ \\
۲ & 20  & 0.111 & $-0.632$ \\
۳ & 30  & 0.222 & $-0.316$ \\
۴ & 40  & 0.333 & \phantom{$-$}0.000 \\
۵ & 100 & 1.000 & \phantom{$-$}1.897 \\
\bottomrule
\end{tabular}
\end{table}

اهمیت مقیاس‌بندی را وقتی بهتر می‌بینیم که چند ویژگی با مقیاس‌های بسیار
متفاوت داشته باشیم. شکل \ref{fig:scaling-effect} دو خوشهٔ مصنوعی از افراد را
بر اساس «سن» (بازهٔ چند ده) و «درآمد» (بازهٔ ده‌ها هزار) نشان می‌دهد. در
نمودار سمت چپ (مقیاس خام)، فاصلهٔ اقلیدسی عملاً فقط به درآمد وابسته است و
سن هیچ نقشی ندارد؛ پس از هر دو نوع مقیاس‌بندی، دو ویژگی وزن قابل‌مقایسه‌ای
پیدا می‌کنند.

\begin{figure}[htbp]
\centering
\includegraphics[width=\linewidth]{figures/ch02_scaling_effect.pdf}
\caption{اثر مقیاس‌بندی بر دو ویژگی با مقیاس متفاوت: از چپ به راست، مقیاس خام،
\LR{Min-Max} و \LR{Z-score}}
\label{fig:scaling-effect}
\end{figure}

پیاده‌سازی این دو تبدیل با \texttt{scikit-learn} به‌صورت زیر است:

\begin{lstlisting}
import numpy as np
from sklearn.preprocessing import MinMaxScaler, StandardScaler

x = np.array([[10], [20], [30], [40], [100]], dtype=float)

x_minmax = MinMaxScaler().fit_transform(x)
x_zscore = StandardScaler().fit_transform(x)

print(x_minmax.ravel().round(3))  # [0.    0.111 0.222 0.333 1.   ]
print(x_zscore.ravel().round(3))  # [-0.949 -0.632 -0.316  0.     1.897]
\end{lstlisting}

\subsection{کاهش داده و کاهش ابعاد}

هدف کاهش داده، رسیدن به بازنمایی کوچک‌تر اما تقریباً برابر از نظر محتوای
اطلاعاتی است. دو رویکرد اصلی عبارت‌اند از:

\begin{itemize}
    \item \textbf{کاهش تعداد نمونه‌ها}: از طریق \dmterm{نمونه‌برداری}{Sampling} یا
          تجمیع داده.
    \item \textbf{کاهش تعداد ویژگی‌ها (کاهش ابعاد)}: از طریق
          \dmterm{انتخاب زیرمجموعه‌ای از ویژگی‌ها}{Feature Selection} یا
          استخراج ویژگی جدید.
\end{itemize}

رایج‌ترین روش استخراج ویژگی، \dmterm{\textbf{تحلیل مؤلفه‌های اصلی}}{Principal Component Analysis (PCA)}
است. ایدهٔ اصلی \LR{PCA} یافتن جهت‌هایی (مؤلفه‌های اصلی) در فضای ویژگی‌هاست که
بیشترین واریانس داده را حفظ می‌کنند؛ با تصویر کردن داده روی چند مؤلفهٔ اول،
می‌توان ابعاد را کاهش داد در حالی که بخش عمدهٔ اطلاعات داده حفظ می‌شود. توجه
شود که \LR{PCA} پیش از اجرا نیازمند استانداردسازی داده است، زیرا بر پایهٔ
واریانس عمل می‌کند.

مراحل کلی \LR{PCA} به‌صورت زیر است:

\begin{enumerate}
    \item استانداردسازی هر ویژگی (میانگین صفر و واریانس یک).
    \item محاسبهٔ ماتریس کوواریانس ویژگی‌ها.
    \item یافتن بردارهای ویژه و مقادیر ویژهٔ این ماتریس؛ بردار ویژه با بزرگ‌ترین
          مقدار ویژه، جهت بیشترین واریانس (مؤلفهٔ اول) است.
    \item انتخاب $k$ مؤلفهٔ اول و تصویر کردن داده روی آن‌ها.
\end{enumerate}

مجموعه دادهٔ \LR{Iris} چهار ویژگی دارد و نمی‌توان آن را مستقیماً روی صفحه
رسم کرد. شکل \ref{fig:pca-iris} نتیجهٔ کاهش آن به دو مؤلفهٔ اصلی را نشان می‌دهد.
دو مؤلفهٔ اول در مجموع حدود ۹۶ درصد واریانس داده را نگه می‌دارند
(۷۳ درصد و ۲۳ درصد)، و با وجود حذف دو بعد، سه گونه هنوز به‌خوبی قابل تشخیص‌اند.

\begin{figure}[htbp]
\centering
\includegraphics[width=0.72\linewidth]{figures/ch02_pca_iris.pdf}
\caption{تصویر مجموعه دادهٔ \LR{Iris} روی دو مؤلفهٔ اصلی (کاهش از چهار بعد به دو بعد)}
\label{fig:pca-iris}
\end{figure}

\begin{lstlisting}
from sklearn.datasets import load_iris
from sklearn.preprocessing import StandardScaler
from sklearn.decomposition import PCA

X = load_iris().data
X_scaled = StandardScaler().fit_transform(X)

pca = PCA(n_components=2)
X_2d = pca.fit_transform(X_scaled)

print(pca.explained_variance_ratio_)  # about [0.730 0.229]
\end{lstlisting}

\subsection{گسسته‌سازی و سلسله‌مراتب مفهومی}
\label{subsec:discretization}

برخی الگوریتم‌ها (مانند نسخه‌های پایهٔ درخت تصمیم یا بیز ساده) بهتر با
داده‌های کیفی کار می‌کنند. \textbf{گسسته‌سازی} فرآیند تبدیل یک ویژگی عددی
پیوسته به تعدادی بازهٔ گسسته (مانند تبدیل سن به گروه‌های «جوان، میان‌سال،
مسن») است. روش‌های رایج شامل بین‌بندی با عرض یکسان\footnote{\LR{Equal-width}}،
بین‌بندی با فراوانی یکسان\footnote{\LR{Equal-depth}} و روش‌های مبتنی بر
خوشه‌بندی هستند.

تفاوت این دو روش بین‌بندی در شکل \ref{fig:binning} دیده می‌شود. در
بین‌بندی با عرض یکسان، طول همهٔ بازه‌ها برابر است اما تعداد رکوردها در آن‌ها
یکسان نیست (وقتی داده چوله است، یک بازه ممکن است تقریباً خالی بماند). در
بین‌بندی با فراوانی یکسان، هر بازه تعداد برابری رکورد دارد اما عرض بازه‌ها
متفاوت است.

\begin{figure}[htbp]
\centering
\includegraphics[width=\linewidth]{figures/ch02_binning.pdf}
\caption{مقایسهٔ بین‌بندی با عرض یکسان (چپ) و بین‌بندی با فراوانی یکسان (راست) روی ویژگی سن}
\label{fig:binning}
\end{figure}

\begin{lstlisting}
import pandas as pd

ages = pd.Series([22, 25, 31, 34, 38, 41, 47, 52, 58, 66])

equal_width = pd.cut(ages, bins=4)    # same interval length
equal_depth = pd.qcut(ages, q=4)      # same number of records

print(equal_width.value_counts().sort_index())
print(equal_depth.value_counts().sort_index())
\end{lstlisting}

\dmterm{\textbf{سلسله‌مراتب مفهومی}}{Concept Hierarchy} نیز به تعریف سطوح مختلف
انتزاع برای یک ویژگی اشاره دارد (مثلاً شهر $\to$ استان $\to$ کشور)، که در
تحلیل‌های چندسطحی کاربرد دارد.

\section{معیارهای شباهت و فاصله}
\label{sec:similarity}

بسیاری از الگوریتم‌های داده‌کاوی (خوشه‌بندی، نزدیک‌ترین همسایه، تشخیص داده پرت)
بر پایهٔ سنجش شباهت یا فاصله میان دو نمونه عمل می‌کنند. رایج‌ترین معیارها
عبارت‌اند از:

\paragraph{فاصلهٔ اقلیدسی}
برای دو بردار $x$ و $y$ در فضای $n$-بعدی:
\[
d(x, y) = \sqrt{\sum_{i=1}^{n} (x_i - y_i)^2}
\]

\paragraph{فاصلهٔ منهتن}
\[
d(x, y) = \sum_{i=1}^{n} |x_i - y_i|
\]

تفاوت این دو فاصله را می‌توان روی صفحه دید. دو نقطهٔ $A=(0,0)$ و $B=(4,3)$ را
در نظر بگیرید. فاصلهٔ اقلیدسی طول پاره‌خط مستقیم میان آن‌هاست:
$\sqrt{4^2+3^2}=5$. فاصلهٔ منهتن مانند حرکت در خیابان‌های شبکه‌ای یک شهر است
که فقط در امتداد افقی و عمودی می‌توان رفت: $4+3=7$ (شکل \ref{fig:distances}).

\begin{figure}[htbp]
\centering
\begin{tikzpicture}[scale=0.95]
\draw[step=1, gray!30, very thin] (-0.4,-0.4) grid (5.4,4.4);
\draw[->, gray] (-0.4,0) -- (5.4,0);
\draw[->, gray] (0,-0.4) -- (0,4.4);
\coordinate (A) at (0,0);
\coordinate (B) at (4,3);
\draw[very thick, blue!70] (A) -- (B);
\draw[very thick, red!70, dashed]
    (0,0) -- (1,0) -- (1,1) -- (2,1) -- (2,2) -- (3,2) -- (3,3) -- (4,3);
\fill[black] (A) circle (2.5pt) node[below left] {$A$};
\fill[black] (B) circle (2.5pt) node[above right] {$B$};
\node[blue!70, font=\small] at (1.0,2.5) {$d_2 = 5$};
\node[red!70, font=\small] at (3.7,1.4) {$d_1 = 7$};

\draw[very thick, blue!70] (6.3,3.6) -- (7.0,3.6);
\node[right, font=\small] at (7.1,3.6) {فاصلهٔ اقلیدسی};
\draw[very thick, red!70, dashed] (6.3,2.9) -- (7.0,2.9);
\node[right, font=\small] at (7.1,2.9) {فاصلهٔ منهتن};
\end{tikzpicture}
\caption{مقایسهٔ فاصلهٔ اقلیدسی و فاصلهٔ منهتن میان دو نقطه در صفحه}
\label{fig:distances}
\end{figure}

\paragraph{فاصلهٔ مینکوفسکی}
تعمیمی از دو فاصلهٔ بالا با پارامتر $p$:
\[
d(x, y) = \left( \sum_{i=1}^{n} |x_i - y_i|^p \right)^{1/p}
\]
که برای $p=1$ به فاصلهٔ منهتن و برای $p=2$ به فاصلهٔ اقلیدسی تبدیل می‌شود.

\paragraph{شباهت کسینوسی}
برای داده‌های با ابعاد بالا و پراکنده (مانند بردارهای متنی)، شباهت کسینوسی
که زاویهٔ میان دو بردار را می‌سنجد، اغلب مناسب‌تر از فاصلهٔ اقلیدسی است:
\[
\text{sim}(x, y) = \frac{x \cdot y}{\|x\| \, \|y\|}
\]

برای ویژگی‌های اسمی یا دودویی نیز معیارهایی مانند
\dmterm{\textbf{ضریب تطابق ساده}}{Simple Matching Coefficient} و
\dmterm{\textbf{ضریب ژاکارد}}{Jaccard Coefficient} به کار می‌روند که در فصل
مربوط به خوشه‌بندی (فصل \ref{ch:clustering}) دوباره به آن‌ها پرداخته می‌شود.
جدول \ref{tab:distance-summary} چهار معیار اصلی این بخش را کنار هم قرار می‌دهد.

\begin{table}[htbp]
\centering
\caption{خلاصهٔ معیارهای فاصله و شباهت}
\label{tab:distance-summary}
\begin{tabular}{@{}p{2.6cm}p{3.9cm}p{5.6cm}@{}}
\toprule
\textbf{معیار} & \textbf{فرمول} & \textbf{کاربرد مناسب} \\
\midrule
اقلیدسی & $\sqrt{\sum_i (x_i-y_i)^2}$ & داده‌های عددی با ابعاد کم و مقیاس یکسان \\
\addlinespace
منهتن & $\sum_i |x_i-y_i|$ & داده‌هایی که مقادیر پرت دارند یا ابعاد آن‌ها مستقل‌اند \\
\addlinespace
مینکوفسکی & $\left(\sum_i |x_i-y_i|^p\right)^{1/p}$ & تعمیم دو معیار بالا؛ پارامتر $p$ قابل تنظیم \\
\addlinespace
کسینوسی & $\dfrac{x\cdot y}{\|x\|\,\|y\|}$ & داده‌های پراکنده و پرابعاد مانند متن \\
\bottomrule
\end{tabular}
\end{table}

محاسبهٔ این معیارها با \texttt{scipy} کوتاه است. برای دو نقطهٔ شکل
\ref{fig:distances}:

\begin{lstlisting}
from scipy.spatial import distance

A, B = [0, 0], [4, 3]

print(distance.euclidean(A, B))   # 5.0
print(distance.cityblock(A, B))   # 7   (Manhattan)
print(distance.minkowski(A, B, p=3))

# cosine *distance* = 1 - cosine similarity
print(1 - distance.cosine([1, 2, 3], [2, 4, 6]))  # 1.0
\end{lstlisting}

\section*{پیاده‌سازی در پایتون}

نمونه‌کدهای این فصل در متن هر بخش آمده‌اند. مجموعهٔ کامل آن‌ها به‌صورت یک
اسکریپت واحد (جایگذاری، مقیاس‌بندی، \LR{PCA}، بین‌بندی و محاسبهٔ فاصله) در بخش
«فصل ۲» پیوست الف (\ref{app:python}) گردآوری شده است.

\section*{خلاصهٔ فصل}

در این فصل با انواع ویژگی‌ها، منابع اصلی مشکلات کیفیت داده، و مراحل کلیدی
پیش‌پردازش شامل پاک‌سازی، یکپارچه‌سازی، تبدیل، کاهش داده و گسسته‌سازی آشنا
شدیم. همچنین رایج‌ترین معیارهای فاصله و شباهت را که در فصل‌های بعدی
(به‌ویژه طبقه‌بندی و خوشه‌بندی) بارها استفاده خواهند شد، مرور کردیم. در
فصل بعد، به اولین وظیفهٔ اصلی داده‌کاوی یعنی طبقه‌بندی می‌پردازیم.

\section*{تمرین‌ها}

\begin{enumerate}
    \item برای هر یک از ویژگی‌های «شمارهٔ ملی»، «دمای بدن»، «رتبهٔ مسابقه»
          و «تعداد فرزندان» نوع آن (اسمی/ترتیبی/بازه‌ای/نسبتی) را مشخص کنید.
    \item فرض کنید در یک مجموعه داده، ۱۵ درصد مقادیر ستون «درآمد» گمشده است.
          با توجه به جدول \ref{tab:missing-strategies}، دو راهکار پیشنهاد دهید و
          مزایا و معایب هرکدام را توضیح دهید.
    \item مقادیر $5, 15, 25, 35, 500$ را هم با \LR{Min-Max} و هم با \LR{Z-score}
          تبدیل کنید و نتیجه را با جدول \ref{tab:scaling-example} مقایسه کنید. کدام روش
          در برابر مقدار ۵۰۰ پایدارتر است؟
    \item فاصلهٔ اقلیدسی و فاصلهٔ منهتن را برای دو بردار
          $x = (2, 4, 6)$ و $y = (5, 1, 8)$ محاسبه کنید و پاسخ را با کد
          \texttt{scipy} بررسی کنید.
    \item چرا شباهت کسینوسی برای مقایسهٔ اسناد متنی معمولاً بهتر از فاصلهٔ
          اقلیدسی عمل می‌کند؟
    \item در شکل \ref{fig:pca-iris}، اگر فقط مؤلفهٔ اول را نگه می‌داشتیم، تفکیک کدام
          دو گونه دشوارتر می‌شد؟ چرا؟
    \item با \texttt{pd.cut} و \texttt{pd.qcut} یک ستون سن دلخواه را بین‌بندی کنید
          و تفاوت تعداد رکوردها در بازه‌ها را با شکل \ref{fig:binning} مقایسه کنید.
\end{enumerate}
